\section{Результаты на ARC-AGI и иллюстративные графики}\label{sec:arc_results}

\DTLloadrawdb{arcbase}{tables/arc_stopping_base.csv}%
\DTLloadrawdb{arcft}{tables/arc_stopping_finetune_compare.csv}%
\DTLloadrawdb{tauhist}{tables/tau_star_hist_finite_discrete_val.csv}%
\newcommand{\TauHistRows}{}%
\DTLforeach*{tauhist}{\ZcA=step,\ZcB=count,\ZcC=probability}{%
  \csname protected@edef\endcsname\DTtabrow{\ZcA & \ZcB & \noexpand\num{\ZcC}\noexpand\\}%
  \expandafter\appto\expandafter\TauHistRows\expandafter{\DTtabrow}%
}%
\newcommand{\ArcBaseRows}{}%
\DTLforeach*{arcbase}{\ZcA=family,\ZcB=strategy,\ZcC=threshold,\ZcD=meantau,\ZcE=regret,\ZcF=tokbest,\ZcG=nlltau}{%
  \csname protected@edef\endcsname\DTtabrow{\noexpand\texttt{\ZcA} & \noexpand\texttt{\ZcB} & \noexpand\num{\ZcC} & \noexpand\num{\ZcD} & \noexpand\num{\ZcE} & \noexpand\num{\ZcF} & \noexpand\num{\ZcG}\noexpand\\}%
  \expandafter\appto\expandafter\ArcBaseRows\expandafter{\DTtabrow}%
}%
\newcommand{\ArcFtRows}{}%
\DTLforeach*{arcft}{\ZcA=family,\ZcB=policybase,\ZcC=policyft,\ZcD=tokbase,\ZcE=tokft,\ZcF=taubase,\ZcG=tauft}{%
  \csname protected@edef\endcsname\DTtabrow{\noexpand\texttt{\ZcA} & \ZcB & \ZcC & \noexpand\num{\ZcD} & \noexpand\num{\ZcE} & \noexpand\num{\ZcF} & \noexpand\num{\ZcG}\noexpand\\}%
  \expandafter\appto\expandafter\ArcFtRows\expandafter{\DTtabrow}%
}%

\subsection{Эмпирическое распределение оракульного \(\tau^\star\)}

На рис.~\ref{fig:arc_tau_hist} показана гистограмма оракульного времени остановки на валидационном подмножестве ARC-AGI для конфигурации \texttt{finite\_discrete} (максимум \(K=8\) внутренних шагов, \(400\) задач в оценке по JSON-артефакту построения графика). Около \(82\%\) массы сосредоточено в \(\tau^\star=0\) в смысле выбранного в коде критерия \texttt{argmax\_masked\_token\_accuracy\_per\_supervision\_step}; заметная доля приходится на \(\tau^\star\in\{1,2,3\}\), хвост до \(K\) ненулевой. Это подтверждает, что адаптивная остановка потенциально содержательна: однородный запуск на полную глубину смешивает простые и сложные задачи.

\begin{figure}[H]
\centering
\includegraphics[width=0.88\linewidth]{arc_tau_star_hist_val.png}
\caption{Эмпирическое распределение оракульного \(\tau^\star\) на валидации (ARC-AGI, ветка \texttt{finite\_discrete}).}
\label{fig:arc_tau_hist}
\end{figure}

Численные счётчики и доли по шагам (источник: \texttt{tables/tau\_star\_hist\_finite\_discrete\_val.csv}, согласован с \texttt{tau\_star\_hist\_val.json}) приведены в табл.~\ref{tab:tau_hist_csv}.

\begin{table}[H]
\centering
\caption{Эмпирическое распределение \(\tau^\star\) по шагам (валидация, \texttt{finite\_discrete})}
\label{tab:tau_hist_csv}
\small
\begin{tabular}{r r S[table-format=1.4]}
\toprule
Шаг & Число задач & {Доля} \\
\midrule
\TauHistRows
\bottomrule
\end{tabular}
\end{table}

\subsection{Сводка по семействам распределений (базовое обучение)}

Табл.~\ref{tab:arc_base} агрегирует лучшую по сетке строку для каждого обученного семейства (метрика отбора: \texttt{token\_acc\_best\_of\_policies}). Полный перебор порогов и стратегий хранится в \texttt{stopping\_eval.json} каждого прогона; табличные значения вынесены в \texttt{tables/arc\_stopping\_base.csv}.

\begin{table}[H]
\centering
\caption{ARC-AGI, валидация: лучшая комбинация правила остановки после базового обучения}
\label{tab:arc_base}
\small
\begin{tabular}{l l S[table-format=1.1] S[table-format=1.2] S[table-format=1.5] S[table-format=1.3] S[table-format=1.3]}
\toprule
Семейство & Стратегия & {Порог} & {\(\overline{\tau}\)} & {Регрет} & {\(\mathrm{tok}_{\mathrm{best}}\)} & {\(\mathrm{NLL}_\tau\)} \\
\midrule
\ArcBaseRows
\bottomrule
\end{tabular}
\end{table}

Здесь \(\overline{\tau}=\) \texttt{best\_mean\_steps} --- среднее число внутренних шагов при выбранном правиле; средний регрет --- \texttt{best\_mean\_regret} из сводки; \(\mathrm{tok}_{\mathrm{best}}=\) \texttt{best\_token\_acc\_best\_of\_policies}; \(\mathrm{NLL}_\tau=\) \texttt{best\_nll}. Доля полностью точных решеток на уровне токенов близка к нулю, что ожидаемо для сложной задачи при пошаговой разметке.

\subsection{Влияние дообучения только оракула}

Табл.~\ref{tab:arc_ft} сопоставляет ту же метрику \(\mathrm{tok}_{\mathrm{best}}\) и среднее число шагов до и после \texttt{oracle\_finetune}. Пороги и стратегии переоптимизируются заново на валидации после дообучения. Данные: \texttt{tables/arc\_stopping\_finetune\_compare.csv}.

\begin{table}[H]
\centering
\caption{Сравнение базового чекпойнта и \texttt{oracle\_finetune} (ARC-AGI)}
\label{tab:arc_ft}
\small
\begin{tabular}{l c c S S S S}
\toprule
 & \multicolumn{2}{c}{Стратегия / порог} & \multicolumn{2}{c}{\(\mathrm{tok}_{\mathrm{best}}\)} & \multicolumn{2}{c}{\(\overline{\tau}\)} \\
\cmidrule(lr){2-3}\cmidrule(lr){4-5}\cmidrule(lr){6-7}
Семейство & база & FT & {база} & {FT} & {база} & {FT} \\
\midrule
\ArcFtRows
\bottomrule
\end{tabular}
\end{table}

Для \texttt{smoothed\_loss} дообучение не меняет оптимальную строку сводки. Для \texttt{finite\_discrete}, \texttt{mixture\_geometric} и \texttt{power} оптимальное правило смещается в сторону более агрессивной ранней остановки (\texttt{hazard} или меньший порог \texttt{cumulative\_probability}), при этом \(\mathrm{tok}_{\mathrm{best}}\) меняется незначительно. Для \texttt{negative\_binomial} наблюдается небольшое снижение \(\mathrm{tok}_{\mathrm{best}}\) после дообучения; интерпретация этого факта обсуждается в разд.~\ref{sec:discussion}.

\subsection{Кривая качество--стоимость и калибровка}

Полные кривые «качество--среднее число шагов'' строятся по всем строкам \texttt{stopping\_eval.json} и здесь не приводятся. На рис.~\ref{fig:quality_cost} оставлен лишь шаблон осей: фактические точки получаются при экспорте JSON в скрипт визуализации.

\begin{figure}[H]
\centering
\fbox{\begin{minipage}{0.82\textwidth}
\centering
\vspace{8mm}
\small ось X: среднее число внутренних шагов; ось Y: токеновая точность / потеря; семейства распределений дают разные траектории при изменении порога.
\vspace{8mm}
\end{minipage}}
\caption{Шаблон кривой качество--стоимость (заполнить по \texttt{stopping\_eval.json})}
\label{fig:quality_cost}
\end{figure}

Reliability-диаграммы для событий \(\{\tau^\star\le k\}\) могут быть построены по полям \texttt{brier}, \texttt{ece} из лучших строк сетки; на рис.~\ref{fig:calibration} --- заготовка.

\begin{figure}[H]
\centering
\fbox{\begin{minipage}{0.82\textwidth}
\centering
\vspace{8mm}
\small ось X: предсказанная вероятность; ось Y: эмпирическая частота; диагональ --- идеальная калибровка \cite{guo2017calibration}.
\vspace{8mm}
\end{minipage}}
\caption{Шаблон калибровочной диаграммы для \(\tau^\star\le k\)}
\label{fig:calibration}
\end{figure}

\subsection{Итоговая таблица для языкового трека (плейсхолдер)}

Табл.~\ref{tab:main_results_template} остается для будущего заполнения при подключении языковых прогонов к тому же коду оценки остановки.

\begin{table}[H]
\centering
\caption{Иллюстративный шаблон итогового сравнения на языковой задаче (не измерено в данной работе)}
\label{tab:main_results_template}
\begin{tabularx}{\textwidth}{lYYYYY}
\toprule
Метод & Потеря речи & Доля совпадений & \(\mathrm{NLL}_{\tau}\) & Среднее число шагов & Экономия \\
\midrule
Фиксированная глубина & \placeholder & \placeholder & не применяется & \placeholder & \placeholder \\
Полный запуск рекурсии & \placeholder & \placeholder & не применяется & \placeholder & \placeholder \\
Эвристика по энтропии & \placeholder & \placeholder & не применяется & \placeholder & \placeholder \\
Конечное распределение & \placeholder & \placeholder & \placeholder & \placeholder & \placeholder \\
Сглаженная разметка & \placeholder & \placeholder & \placeholder & \placeholder & \placeholder \\
Смесь геометрических & \placeholder & \placeholder & \placeholder & \placeholder & \placeholder \\
Отрицательное биномиальное & \placeholder & \placeholder & \placeholder & \placeholder & \placeholder \\
\bottomrule
\end{tabularx}
\end{table}
